\subsection*{Erd\H{o}s problem \#363}

\noindent\textbf{1) FORMAL RESTATEMENT.}
The finiteness proposal asks whether disjoint positive-integer intervals, each of length at least four, can have a square total product only finitely often. We disprove it even with exactly three intervals, all of length four.

\noindent\textbf{2) CONTEXT AND METHOD.}
Ulas first disproved the proposal using four blocks. We reconstruct the stronger three-block Pell family of Bauer--Bennett, Section 4 of \emph{On a question of Erd\H{o}s and Graham}. The argument is explicit and elementary, with no novelty claim.

\noindent\textbf{3) AN INFINITE PELL FAMILY WITH THE REQUIRED CONGRUENCES.}
Start with $(u_0,v_0)=(265,153)$ and iterate
\[
 u_{t+1}=97u_t+168v_t,\qquad
 v_{t+1}=56u_t+97v_t.
\]
Direct calculation gives $265^2-3\cdot153^2=-2$ and $97^2-3\cdot56^2=1$. Multiplication of $u+v\sqrt3$ by $97+56\sqrt3$ therefore preserves
\[
 u^2-3v^2=-2.
\]
Both coordinates increase strictly. Modulo $4$, the recurrence preserves $u\equiv1$; modulo $2$, it preserves $v\equiv1$. Thus every pair defines positive integers
\[
 x_1=(u-5)/4,\qquad x_2=(v-3)/2,\qquad x_3=(u+1)/2.
\]
Let $I_i=\{x_i,x_i+1,x_i+2,x_i+3\}$.

\noindent\textbf{4) DISJOINTNESS.}
Since $u\ge265$ and $v^2=(u^2+2)/3$, we have
\[
 4u/7<v<3u/5.
\]
The lower inequality follows from $1/3>16/49$, and the upper from $u^2>25$. Consequently
\[
 x_2-x_1=(2v-u-1)/4>(u/7-1)/4>4,
\]
\[
 x_3-x_2=(u-v+4)/2>(2u/5+4)/2>4.
\]
The three intervals are disjoint, and $x_1\to\infty$ along the recurrence, giving genuinely different collections.

\noindent\textbf{5) EXACT SQUARE IDENTITY.}
Group the twelve factors as follows:
\[
 x_1(x_3+2)x_2(x_2+3)
 =\frac{(u^2-25)(v^2-9)}{32}
 =\frac{(u^2-25)^2}{96},
\]
\[
 (x_1+1)x_3(x_2+1)(x_2+2)
 =\frac{(u^2-1)(v^2-1)}{32}
 =\frac{(u^2-1)^2}{96},
\]
\[
 (x_1+2)(x_1+3)(x_3+1)(x_3+3)
 =\frac{(u+3)^2(u+7)^2}{64}.
\]
Multiplying gives
\[
 \prod_{i=1}^3\prod_{h=0}^3(x_i+h)
 =\left(\frac{(u^2-25)(u^2-1)(u+3)(u+7)}{768}\right)^2.
\]
The right-hand square root is rational, and the left side is an integer; a rational number whose square is an integer is itself an integer, by writing it in lowest terms. Thus the product is an integer square, without any unstated divisibility assumption on the displayed numerator.

For the initial pair, the starts are $65,75,133$. These three length-four blocks already provide a concrete counterexample, while the recurrence supplies infinitely many.

\noindent\textbf{6) VERIFICATION AND FINAL.}
The recurrence proves infinitude and preserves both the Pell identity and the integrality congruences. Every interval is positive, has exactly four elements, and is disjoint from the others. \textbf{SOLVED: the proposed finiteness statement is false.} The proof establishes an explicit infinite family and makes no assertion about the separate two-block case.

\subsection*{Sources}
\url{https://www.erdosproblems.com/363}; \url{https://personal.math.ubc.ca/~bennett/Erdos-Graham.pdf}.

\section{COMPLETION ESTIMATE}
\textbf{COMPLETION: 100\%}
